\section{Vanishing denominators without spectral collapse}\label{vanishingdenominator}
The preceding family has small potential values but a uniformly positive
noise denominator. We next rule out deriving a positive denominator floor
from the original hypotheses. The same construction shows why excluding
individual points with a large ratio is too strong a proposed remedy.

\begin{theorem}\label{vanishingdenominatortheorem}
Fix $0<\rho,b<1$ and $C>0$. There are energy-normalized positive
potentials $\widetilde h_L$ satisfying all of \eqref{admissible} such that,
at the all-plus vertex $x_0$, for a constant $\varepsilon>0$,
\[
 \widetilde h_L(x_0)\sim e^{\varepsilon L},\qquad
 (P\widetilde h_L)(x_0)\sim e^{-\varepsilon L},\qquad
 \widetilde\pi_L(x_0)\longrightarrow1.
\]
Thus $\widetilde U_L(x_0)\sim e^{2\varepsilon L}$, where
$\widetilde U_L=\widetilde h_L/(P\widetilde h_L)$. Nevertheless there
are constants $A<\infty$, $a>0$ such that, for all sufficiently large $L$,
\[
 \norm{\widetilde\pi_L-1}_s+\norm{\widetilde U_L-1}_s
       \leq Ae^{-aL},\qquad
 \max_{1\leq j\leq2^L}
 |\lambda_j(K_{\widetilde h_L})-\lambda_j(P)|\leq Ae^{-aL}.
\]
Eigenvalues are ordered with multiplicities. In particular the full gap
tends to $1-\rho$. Constants and the eventual threshold may depend on
the fixed parameters; no uniformity as $\rho$ approaches its endpoints
is claimed.
\end{theorem}
\begin{proof}
We use the classical Bernoulli concentration bound, including unequal
parameters: if $S$ sums $L$ independent Bernoulli variables, then
\[
 \Pr(S-\E S\geq tL),\ \Pr(S-\E S\leq-tL)
       \leq e^{-2t^2L}\qquad(t>0).
\]
Here is the needed proof. The log moment-generating function of a
centered Bernoulli variable has value and first derivative zero at zero;
its second derivative is the variance of a tilted Bernoulli law, at most
$1/4$. Integrating twice bounds it by $\lambda^2/8$ for every real
$\lambda$. Independence and exponential Markov at $\lambda=\pm4t$
give the two bounds. No equal-parameter assumption is used.

Put
\[
 p=\frac{1-\rho}{2},\quad q=p+\rho p=\frac{1-\rho^2}{2},\quad
 a_0=\frac{1+b}{2},\quad
 \delta=\min\{\rho p/16,\rho^2/16\},\quad r_0=2\delta/\rho.
\]
Let $d(x)$ be the Hamming distance from $x_0$ and define
\begin{align*}
 B_L&=\{x:|d(x)/L-p|\leq\delta\},\\
 A_L&=\{x:|d(x)/L-q|\leq2\delta\},\\
 D_L&=\{x:(P1_{B_L})(x)>1-a_0\}.
\end{align*}
In this proof $D_L$ is a set. We have $r_0\leq p/8$,
$q-p\geq16\delta$ and $1/2-q\geq8\delta$. Thus $A_L$ and $B_L$
are disjoint. Conditional on $x$, the next noise step has independent
coordinate flips and
\[
 \E_x d(Y)=pL+\rho d(x).
\]
For $x\in B_L$ this mean lies between $(q-\rho\delta)L$ and
$(q+\rho\delta)L$, so
\[
 P(x,A_L)\geq1-2e^{-2\delta^2L},\qquad
 P(x_0,B_L)\geq1-2e^{-2\delta^2L}.
\]
If $d(x)/L\geq r_0$, the mean is at least $(p+2\delta)L$ and the
upper boundary of $B_L$ is $(p+\delta)L$. Therefore
$P(x,B_L)\leq e^{-2\delta^2L}$. For all sufficiently large $L$,
\[
 x_0\in D_L\subset\{x:d(x)/L<r_0\}.
\]
The last ball is disjoint from both shells, since $r_0<p-\delta$.

Choose the strictly positive constant
\[
 c=\min\left\{
 \begin{gathered}
 2\delta^2,\ 2(p-r_0)^2,\ 2(q-2\delta-p)^2,\\
 2(1/2-q-2\delta)^2,\ 2(1/2-p-\delta)^2,\ 2(1/2-r_0)^2
 \end{gathered}\right\}.
\]
Under uniform measure $d(x)$ is binomial with mean $L/2$, whereas
under $P(x_0,\cdot)$ it is binomial with mean $pL$. The concentration
bounds and the containment of $D_L$ consequently give
\begin{align}
 \mu(A_L),\mu(B_L),\mu(D_L)&\leq e^{-cL},\nonumber\\
 P(x_0,A_L),P(x_0,D_L)&\leq e^{-cL},\nonumber\\
 P(x_0,B_L^c)&\leq2e^{-cL},\qquad
 \inf_{x\in B_L}P(x,A_L)\geq1-2e^{-cL}.
 \label{denominatorshellbounds}
\end{align}
These estimates include shell boundaries and the return to $x_0$.

Set $\varepsilon=c/(16s)$, $M=e^{\varepsilon L}$, and $E_L=A_L\cup D_L$.
Define the actual potential and its product by
\[
 H=\begin{cases}M&\text{on }E_L,\\M^{-1}&\text{on }B_L,\\1&\text{elsewhere},\end{cases}
 \qquad G=PH,\quad R=HG,\quad Z=\E R,\quad h=H/\sqrt Z.
\]
Always $M^{-1}\leq H,G\leq M$. On $E_L$, $R\geq1$. On $B_L$,
$G\geq MP1_{A_L}$, hence $R\geq1-2e^{-cL}\geq a_0$ eventually.
Outside $E_L\cup B_L$, one has $H=1$ and $x\notin D_L$, whence
$G(x)\geq1-P(x,B_L)\geq a_0$. Thus $R\geq a_0$ at every vertex.

At $x_0\in D_L$, splitting the noise expectation over the three
regions gives
\[
 G(x_0)=M^{-1}P(x_0,B_L)+P(x_0,(B_L\cup E_L)^c)+MP(x_0,E_L).
\]
Consequently
\[
 P(x_0,B_L)\leq MG(x_0)
       \leq1+2Me^{-cL}+2M^2e^{-cL}\longrightarrow1,
\]
because $2\varepsilon<c$. Hence $G(x_0)\sim M^{-1}$,
$R(x_0)\to1$, and $H(x_0)/G(x_0)\sim M^2$.

For finite $r\geq1$ write $d_{r,L}=\norm{H-1}_r$. The support
bounds in \eqref{denominatorshellbounds} imply
\[
 d_{r,L}\leq M(2e^{-cL})^{1/r}+e^{-cL/r},\qquad
 \norm{G-1}_r\leq d_{r,L}.
\]
The second inequality is noise contraction. H\"older therefore gives
\[
 \norm{R-1}_s\leq d_{2s,L}(1+d_{2s,L})+d_{2s,L}.
\]
Since $\varepsilon<c/(2s)$, this error decays exponentially, as does
$|Z-1|$. Thus $\pi=R/Z$ has exponentially small $\norm{\pi-1}_s$,
$\E\pi^s\to1$, and eventually $\pi\geq b_1:=(a_0+b)/2>b$.
Normalization preserves $h(x_0)\sim M$, $(Ph)(x_0)\sim M^{-1}$
and $\pi(x_0)\to1$.

Crucially, for $U=H/G$, use the proved dimension-dependent bound
$G\geq M^{-1}$, not a fixed denominator floor. It gives
\begin{align}
 \norm{U-1}_s&\leq M\norm{H-G}_s\leq2M d_{s,L}\nonumber\\
 &\leq2\,2^{1/s}e^{-(c/s-2\varepsilon)L}
          +2e^{-(c/s-\varepsilon)L}.\label{vanishingdenominatorratio}
\end{align}
Both rates are positive. Lemma~\ref{criticalsandwich} controls the
entire actual spectrum by this ratio error.

It remains to impose exact matching. Subtract the finite Walsh sum
\[
 e_L=-\sum_{1\leq|S|\leq k_L}(\E\pi\chi_S)\chi_S,
 \qquad \widetilde\pi=\pi+e_L.
\]
Each coefficient is at most $\norm{\pi-1}_1$, and for all large $L$
there are at most $k_LL^{k_L}=\exp(o(L))$ terms. Therefore
$\xi_L=\norm{e_L}_\infty$ is exponentially small after decreasing
the positive rate. Orthogonality gives exact mean one and all requested
zero coefficients. The path $\pi+\tau e_L$ has floor $b$ eventually,
and $\norm{\widetilde\pi-1}_s$ remains exponentially small, giving
the eventual cap two.

Lemma~\ref{inversestationary} reconstructs the unique positive
$\widetilde h$ with $\widetilde hP\widetilde h=\widetilde\pi$ and gives
\[
 \eta_L:=\norm{\log(\widetilde h/h)}_\infty
 \leq\frac{\lceil(L+\log_2(1/b))/2\rceil+2}{2b}\xi_L.
\]
This is exponentially small. Positivity of $P$ transfers the same
multiplicative bounds to $P\widetilde h$ and $Ph$, preserving the
pointwise asymptotics, including the exponentially small denominator.
Also $e^{-2\eta_L}U\leq\widetilde U\leq e^{2\eta_L}U$, so
\[
 \norm{\widetilde U-1}_s
 \leq e^{2\eta_L}\norm{U-1}_s+e^{2\eta_L}-1.
\]
Combining this with \eqref{vanishingdenominatorratio} and
Lemma~\ref{criticalsandwich} proves the claimed norm and full-spectrum
bounds. Polynomial factors and the subexponential Walsh count are absorbed
by a smaller positive exponential rate. Fill all finitely many early
dimensions with $\widetilde h=1$, preserving every original constraint.
\end{proof}

Theorem~\ref{vanishingdenominatortheorem} refutes a dimension-uniform
positive $Ph$ floor and a uniform pointwise ceiling for $h/Ph$, even
under exact matching. It does not refute an aggregate killed-return or
cut-flow estimate. Small denominators and large ratios at individual
vertices coexist here with a full spectrum approaching Boolean noise.
The original gap question remains open; no original hypothesis was removed.
